
\section{Executable Research Objects}\fbeg{}\begin{frame}{\ft{Executable Research Objects}}


\vspace{-.5em}	

{\Large\fontfamily{uhv}\selectfont

\hspace*{20pt}\begin{minipage}{.96\textwidth}
\vspace{4pt}
		
\begin{lightquadblockc}{0,0.1,0.1}{\parbox{21cm}{\centering \vspace{3pt}}}
\begin{center}\begin{minipage}{\textwidth}
{\Huge \setlength{\leftmargini}{30pt}\begin{enumerate}
\dmitem {\dsp} A way to achieve FAIR data transparency: data, code, and multimedia bundled togethers
\vspace{20pt}
\dmitem {\dsp} Research Object \q{Bundles} follow a constrained folder layout and required JSON manifest, 
with other recommended metadata
\vspace{20pt}
\dmitem {\dsp} Executable Research Objects are less formally defined, but with an emphasis on code sharing
\vspace{20pt}
%\footnotemark[2]\footnotetext{\Large \raisebox{7pt}{\em{2}} 
%	\hspace{7pt}See Slides 11 and 12 for details.}
\end{enumerate}
}\end{minipage}
\end{center}
\end{lightquadblockc}
\end{minipage}

}

\vspace{-11pt}
{\thrule}
\hspace{32pt}\begin{minipage}{.8\textwidth}
{\fontsize{22}{28}\selectfont \setlength{\leftmargini}{3pt}
Other related terms are \q{Research Compendia},
\q{Story Maps} (in \resizebox{!}{14pt}{ArcGIS}), \q{Data Infrastructure Building Blocks} (\resizebox{!}{14pt}{DIBBs}),
\q{Open Data Cube} (\resizebox{!}{14pt}{ODC}), or 
\q{The Whole Tale} (a \resizebox{!}{14pt}{DIBBS} initiative)

}
\end{minipage}

\end{frame}

